\section{Introduction}

LLM inference has become a persistent workload in cloud and high-performance
computing (HPC) environments.  Unlike training, inference is paid repeatedly:
every prompt triggers model execution, KV-cache traffic, communication, and a
period of node occupancy.  Its energy cost is highly scenario dependent.  A
long-context scientific assistant emphasizes compute-intensive prefill; a
coding assistant with long completions emphasizes memory-intensive decode; a
70B or 405B model introduces tensor or pipeline parallelism; and a multi-node
deployment adds network transfer, synchronization, and idle gaps.  Even
configurations with similar throughput can therefore have different joules per
token and different time-to-first-token (TTFT) or time-per-output-token (TPOT).

Recent measurement efforts make this cost observable.  From Words to Watts
quantified multi-GPU LLM inference energy \cite{samsi2023words}; studies of
inference engines and workload geometry expose large software- and
phase-dependent differences \cite{niu2025engines,fernandez2025energy}; and
\bench\ provides declarative, phase-aligned measurement from a single GPU to
multi-node models as large as Llama-3.1-405B \cite{niu2026tokenpowerbench}.
These systems answer an essential question: \emph{what did this inference run
consume?}

Operators must answer a different question before committing scarce GPU time:
\emph{which configurations are worth running?}  The cross product of prompt
and output lengths, batch limits, engines, precisions, tensor and pipeline
parallelism, placement, and power settings easily creates thousands of
candidates.  Exhaustive measurement is particularly costly for large models,
where one candidate may already require an entire node or several nodes.
Conversely, a simulator calibrated on one GPU cannot faithfully invent
NVLink/NVSwitch collectives, inter-node fabric behavior, pipeline bubbles, or
scheduler contention.  The challenge is therefore \textbf{not only better
energy prediction, but better evidence allocation}: deciding when historical
data are sufficient, when a small sandbox probe can remove uncertainty, and
when a full-scale run is unavoidable.

We propose \system, an agentic multi-fidelity Pareto search framework for
energy-efficient LLM inference on GPU clusters, built on top of \bench.  A user
supplies a model, workload trace or token distribution, target hardware,
serving-engine choices, SLOs, and a measurement budget.  A constrained
language-model planner converts that intent into a typed search task and
chooses an experiment subgoal; a probabilistic policy then selects an action
$a=(x,\ell)$ comprising a serving configuration and evidence fidelity.  The
agent may cheaply simulate an in-domain candidate, calibrate phase behavior on
one GPU, characterize collectives within a node, sample a multi-node scaling
point, or verify a top-ranked configuration at target scale.  Predictions
carry uncertainty and provenance; only accepted telemetry is labeled measured.

The design makes the agent operational rather than decorative.  The LLM does
not estimate joules by free-form reasoning.  Deterministic tools perform
feasibility checks, simulation, regression, telemetry integration, and Pareto
analysis.  The agent interprets the scenario, selects tools and calibration
levels, handles failures, and explains why additional evidence is or is not
worth its GPU-hour cost.  This tool-grounded, validation-driven structure is
well aligned with agentic HPC systems, where autonomous decisions must remain
bounded, inspectable, and reproducible.

This paper makes four contributions:
\begin{itemize}
    \item We formulate \emph{agentic multi-fidelity Pareto search for
    energy-efficient LLM inference}: finding SLO-feasible nondominated
    configurations while minimizing the real GPU-hours required to identify
    and verify them.
    \item We design a phase- and communication-aware \emph{Energy Twin} that
    separates idle, prefill, decode, intra-node communication, inter-node
    communication, and serving overhead, with uncertainty attached to every
    estimate.
    \item We introduce \emph{Intent-Conditioned Pareto Information Gain}, a
    decision-directed acquisition criterion that jointly selects configuration
    and fidelity according to expected information about the SLO-feasible
    high-fidelity Pareto set per GPU-hour.
    \item We realize the method as a bounded hybrid agent with typed subgoals,
    interchangeable replay and cluster executors, failure recovery, provenance,
    and mandatory target-scale verification, and define an evaluation that
    separates model, policy, and agent contributions.
\end{itemize}
